\section*{29. 特征为 $p$ 的有限线性群的不变量有限性定理}

\begin{center}
Nachr. v. d. Ges. d. Wiss. zu Göttingen 1926, S. 28--35
\end{center}

由 R. Courant 在 1926 年 7 月 30 日会议上报告。

我在下面证明，有限线性替换群的不变量的著名有限性定理\footnote{关于一个初等的有限性证明，参见 E. Noether, Der Endlichkeitssatz der Invarianten endlicher Gruppen. Math. Ann. 77 (1916), S. 89--92.} 仍然成立，如果作为系数域所取的不是通常的数域，而是任意域--因而尤其可以是特征为 $p$ 的域。\footnote{关于域论概念，参见 E. Steinitz, Algebraische Theorie der Körper. J. f. M. 137 (1910), S. 167--309。特征见 \S\ 4，根域见 \S\ 12，商域见 \S\ 3。} 不过，一旦群的阶能被 $p$ 整除，已知的有限性证明便失效；必须引入更深的算术事实，也就是下面这个此前只在特征零情形已知的

\textbf{有限性判别准则}\footnote{参见 E. Noether, Algebraische und Differentialinvarianten. Jahresber. d. d. Math.-Ver. 32 (1923), S. 177--184, Nr. 3, Satz 4 以及那里列出的文献。在该判别准则的最后一句中，误写为“关于 $\mathfrak R$ 的模基”（整量的基本系统），而不是“关于 $\mathfrak R$ 的某个子环”。}：一个由多项式--更一般地，由有理函数--构成、系数来自一个域的整环 $\mathfrak S$，当且仅当 $\mathfrak S$ 中含有一个有限子环 $\mathfrak R$，使得 $\mathfrak S$ 在 $\mathfrak R$ 上代数地整依赖时，$\mathfrak S$ 是有限的。于是 $\mathfrak R$ 的每一个整基，都通过 $\mathfrak S$ 关于 $\mathfrak R$ 的某个子环的一个模基，补全为 $\mathfrak S$ 的一个整基。

这个判别准则--它本身也有独立的兴趣--依赖于有理基的存在，以及这样一个事实：从一个域过渡到其有限扩张的整量时，除子链定理仍然保持。第二个定理已由 Artin 和 v. d. Waerden 在前一篇札记中对任意特征证明，不过只是在原始范围满足某种有限性前提时；在有限群情形这个前提是满足的（根环给出一个有限扩张）。对于有理基的存在，我给出一个对任意域作为系数域都有效的证明；由此我便能在上述限制下证明有限性判别准则。

对于有限群的不变量，有限性判别准则之所以满足，依据的是对称函数原理。然而，在特征零情形，Galois 预解式的系数已经给出整基；而在一般情形，它们只生成有限性判别准则所要求的有限子环。由这个有限子环的存在，按同样的推理可推出“相对”不变量和“模”不变量的有限性。

这里考察的不变量，作为特殊情形，包括 Dickson 及其学派\footnote{参见 Dickson, On Invariants and the theory of Numbers. The Madison Colloquium 1913。较新的文献见此后出版的 Transactions of the American Mathematical Society 各卷。} 所处理的“模 $p$ 的射影不变量”。在这里，有限性定理此前虽已对模不变量已知，却尚未对一般的“形式”不变量已知。

\subsection*{\S\ 1. 有限性判别准则}

1. \textbf{有理基存在定理。} 第一种表述：由 $n$ 个不定元的有理函数组成、系数来自一个域 $P$ 的每个系统 $S$，都有一个有限有理基；也就是说，系统中存在有限多个函数，使得系统中的每个函数都能用这有限多个函数、以 $P$ 中系数作有理表示。

第二种表述：设 $\mathfrak K$ 是 $P$ 与 $P(x_1,\ldots,x_n)$ 之间的一个中间域--其中 $x_1,\ldots,x_n$ 表示不定元--且超越次数为 $t\leq n$；设 $y_1,\ldots,y_t$ 是 $\mathfrak K$ 中的一个不可约系统\footnote{按 Steinitz 的说法，一个系统称为不可约，是指它由代数无关函数组成。一个系统称为超越次数为 $t$，是指它至少含有一个由 $t$ 个元素组成的不可约系统，但不含有更多元素的不可约系统。关于下面使用的不可约系统的简单定理，参见 Steinitz, \S\ 22。}。那么 $\mathfrak K$ 是 $P(y_1,\ldots,y_t)$ 的有限代数扩张。

这两种表述事实上等价。只需对所有中间域证明有理基的存在即可；因为由一个系统 $S$ 导出的域成为这样一个中间域 $\mathfrak K$，而 $\mathfrak K$ 的有理基立即给出 $S$ 的有理基，这是因为 $\mathfrak K$ 的每个元素都成为 $S$ 中有限多个函数的有理组合。另一方面，由第二种表述可知，$\mathfrak K$ 中任意一个不可约系统 $y_1,\ldots,y_t$，再加上表征 $\mathfrak K$ 作为 $P(y_1,\ldots,y_t)$ 有限扩张的有限多个函数，就给出这样一个有理基。反过来，如果按第一种表述给出了 $\mathfrak K$ 的一个有理基，那么按超越次数的定义，每个基元素都必定代数依赖于 $P(y_1,\ldots,y_t)$；于是 $\mathfrak K$ 就由此被表征为 $P(y_1,\ldots,y_t)$ 的有限代数扩张。

下面在第二种表述下证明该定理。先设 $\mathfrak K$ 与 $P(x_1,\ldots,x_n)$ 具有相同的超越次数 $n$；设 $y_1,\ldots,y_n$ 是 $\mathfrak K$ 中的一个不可约系统，因而也是 $P(x_1,\ldots,x_n)$ 中的一个不可约系统。由于 $y_1,\ldots,y_n$ 的超越次数为 $n$，每个 $x_i$ 都在 $P(y_1,\ldots,y_n)$ 上代数；所以 $P(x_1,\ldots,x_n)$，从而其中间域 $\mathfrak K$，成为 $P(y_1,\ldots,y_n)$ 的有限代数扩张。

现在设 $t<n$；将 $x$ 的编号选定，使得 $y_1,\ldots,y_t,x_{t+1},\ldots,x_n$ 构成 $P(x_1,\ldots,x_n)$ 中的一个不可约系统。于是 $x_{t+1},\ldots,x_n$ 关于 $P(y_1,\ldots,y_t)$ 不可约，并且由代数依赖的传递律，也关于 $\mathfrak K$ 以及 $P(y_1,\ldots,y_t)$ 与 $\mathfrak K$ 之间的每个中间域 $\mathfrak M$ 不可约。这就是说，$\overline{\mathfrak K}=\mathfrak K(x_{t+1},\ldots,x_n)$ 以及 $\overline{\mathfrak M}=\mathfrak M(x_{t+1},\ldots,x_n)$，分别成为 $\mathfrak K$ 以及 $\mathfrak M$ 的纯超越扩张；因而 $\overline{\mathfrak K}$ 中每个不属于 $\mathfrak K$ 的元素都关于 $\mathfrak K$ 超越，同样，$\overline{\mathfrak M}$ 中每个不属于 $\mathfrak M$ 的元素都关于 $\mathfrak M$ 超越。再令 $\overline P$ 等于纯超越扩张 $P(x_{t+1},\ldots,x_n)$；于是 $y_1,\ldots,y_t$ 关于 $\overline P$ 不可约。

现在，把 $\overline P$ 作为系数域，就可以把证明归结到上面已经处理过的情形。由
\[
 P<\mathfrak K<P(x_1,\ldots,x_t,x_{t+1},\ldots,x_n)
\]
\footnote{$P<\mathfrak K$ 表示：$P$ 包含在 $\mathfrak K$ 中，即 $P$ 是 $\mathfrak K$ 的子域。} 可推出
\[
 \overline P<\overline{\mathfrak K}<\overline P(x_1,\ldots,x_t),
\]
其中 $\overline{\mathfrak K}$ 关于 $\overline P$ 的超越次数为 $t$，与 $\overline P(x_1,\ldots,x_t)$ 相同。因此，$\overline{\mathfrak K}$ 关于 $\overline P$ 作为系数域，有一个有限有理基，记为 $y_1,\ldots,y_t; z_1,\ldots,z_s$。这里 $z$ 可以选自 $\mathfrak K$，因为 $\overline{\mathfrak K}$ 是 $\mathfrak K$ 与 $\overline P$ 的合成域。还需证明 $\mathfrak K$ 等于 $\mathfrak M=P(y_1,\ldots,y_t;z_1,\ldots,z_s)$。$\mathfrak M$ 是 $P(y_1,\ldots,y_t)$ 与 $\mathfrak K$ 之间的一个中间域；所以 $\mathfrak K$ 的每个元素都关于 $\mathfrak M$ 代数。另一方面，$\overline{\mathfrak M}$ 等于 $\overline{\mathfrak K}$，因而 $\mathfrak K$ 包含在 $\overline{\mathfrak M}$ 中。于是 $\mathfrak K$ 成为 $\mathfrak M$ 与 $\overline{\mathfrak M}$ 之间的中间域；又由于 $\overline{\mathfrak M}$ 中每个不属于 $\mathfrak M$ 的元素都关于 $\mathfrak M$ 超越，所以必然有 $\mathfrak K$ 与 $\mathfrak M$ 相同。

推论：若 $P<\mathfrak K<\mathfrak L<P(x_1,\ldots,x_n)$，并且 $\mathfrak L$ 关于 $\mathfrak K$ 代数，则 $\mathfrak L$ 是 $\mathfrak K$ 的有限代数扩张。因为 $\mathfrak L$ 的一个有理基中的每个元素于是都关于 $\mathfrak K$ 代数；从而 $\mathfrak L$ 被表征为 $\mathfrak K$ 的有限扩张。

2. \textbf{有限性判别准则的证明。} 这个条件显然是必要的；因为如果 $\mathfrak S$ 是有限整环，那么 $\mathfrak S$ 本身就可看作判别准则中出现的有限子环 $\mathfrak R$。

这个条件在如下前提下是充分的：若系数域具有特征 $p$，则根域 $P^{1/p}$ 是该系数域\footnote{见第28页脚注2。} 的有限扩张；而在特征零情形，不需要任何限制性前提。为了证明这一点，需要证明 $\mathfrak S$ 关于 $\mathfrak R$ 的某个子环有一个有限模基。

令 $\mathfrak K$ 表示 $\mathfrak R$ 的商域\footnote{见第28页脚注2。}，令 $\mathfrak L$ 表示 $\mathfrak S$ 的商域。这些商域存在，因为这里所涉及的是无零因子的环；并且有
\[
 P<\mathfrak K<\mathfrak L<P(x_1,\ldots,x_n).
\]
因此，由 1. 下的推论，$\mathfrak L$ 成为 $\mathfrak K$ 的有限代数扩张域；又按假设，$\mathfrak S$ 中每个元素都关于 $\mathfrak R$ 整，于是 $\mathfrak S$ 成为由 $\mathfrak L$ 中所有关于 $\mathfrak R$ 整的元素组成的环 $\mathfrak S'$ 的子环。但是，由于并未假设 $\mathfrak R$ 在 $\mathfrak K$ 中整闭，有限扩张下除子链定理保持这一事实不能直接引用。必须先过渡到 $\mathfrak R$ 的一个同样有限的子环 $\mathfrak T$，使得这种整闭性成立。

先假设系数域 $P$ 含有无穷多个元素。那么可从 $\mathfrak R$ 中选出一个不可约系统 $g_1(x),\ldots,g_t(x)$，使得 $\mathfrak R$，从而 $\mathfrak S$，关于
\[
 \mathfrak T=P[g_1(x),\ldots,g_t(x)]
\]
整。事实上，设 $f_1(x),\ldots,f_k(x)$ 是按假设存在的 $\mathfrak R$ 的整基；若它们不构成不可约系统，则在必要地重新编号之后，设 $F(f_k;f_1,\ldots,f_{k-1})=0$ 是 $f_k$ 所满足的一个关系。若把 $f_1,\ldots,f_{k-1}$ 替换为
\[
 f_1+t_1f_k,\ldots,f_{k-1}+t_{k-1}f_k,
\]
则当 $t_i$ 取为不定元时，$f_k$，从而也包括 $f_1,\ldots,f_{k-1}$，都关于这些新函数整；由于 $P$ 有无穷多个元素，可把 $t_i$ 特化为 $P$ 中的元素，使这种整依赖保持。于是经过有限多步，并利用整依赖的传递律，就得到所求的环 $\mathfrak T=P[g_1(x),\ldots,g_t(x)]$。此时 $\mathfrak L$ 是 $\mathfrak T$ 的商域 $P(g_1(x),\ldots,g_t(x))$ 的有限扩张域，而由 $\mathfrak L$ 中所有关于 $\mathfrak R$ 整的元素组成的环 $\mathfrak S'$，与由 $\mathfrak L$ 中所有关于 $\mathfrak T$ 整的元素组成的环相同。

在 $\mathfrak T$ 中，关于所有理想的系统有除子链定理，并且 $\mathfrak T$ 在其商域中整闭；二者都因为 $\mathfrak T$ 同构于 $t$ 个不定元的多项式域。由 $P^{1/p}$ 关于 $P$ 有限这一假设，还可推出 $\mathfrak T^{1/p}$ 关于 $\mathfrak T$ 有限。\footnote{参见 Hilbert, Über die vollen Invariantensysteme, \S\ 1. Math. Ann. 42 (1893), S. 313--373.}\footnote{参见引言中引用的 Artin-v. d. Waerden 札记。} 因此，在 $\mathfrak S'$ 中，对所有 $\mathfrak T$-模组成的系统，除子链定理成立；特别地，$\mathfrak S$ 作为 $\mathfrak S'$ 中包含 $\mathfrak T$ 的一个子环，有一个有限 $\mathfrak T$-模基。此外，如果 $\mathfrak L$ 是 $P(g_1(x),\ldots,g_t(x))$ 的第一类扩张，则 $\mathfrak S$ 的有限 $\mathfrak T$-模基的存在不再需要对系数域 $P$ 作任何限制性前提；因此特别是在 $P$ 具有特征零时总是成立。\footnote{例如参见 E. Noether, Abstrakter Aufbau der Idealtheorie in algebraischen Zahl- und Funktionenkörpern, \S\ 3. Math. Ann. 96 (1926), S. 26--61.}

设 $h_1(x),\ldots,h_s(x)$ 是 $\mathfrak S$ 的一个 $\mathfrak T$-模基；那么对 $\mathfrak S$ 中任意 $f(x)$ 的表示
\[
 f(x)=A_1(g)h_1(x)+\cdots+A_s(g)h_s(x)
\]
表明 $g_1(x),\ldots,g_t(x),h_1(x),\ldots,h_s(x)$ 构成所求的整基；这里 $A$ 作为 $\mathfrak T$ 的元素，表示 $g$ 中的多项式。

若域 $P$ 只含有有限多个元素，那么通过添加一个不定元 $u$--也就是一个关于 $\mathfrak S$ 超越的元素--所得的域 $P(u)$ 仍含有无穷多个元素。因此，环 $\mathfrak S(u)$ 关于 $P(u)$ 作为系数域有一个有限整基；并且由于 $\mathfrak S(u)$ 是 $\mathfrak S$ 与 $P(u)$ 的合成环，这个基可以通过按 $u$ 的幂分解而从 $\mathfrak S$ 中选出，不过 $\mathfrak S$ 的基元素 $g_i$ 现在将不再构成不可约系统。这样，$\mathfrak S$ 中每个多项式 $f(x)$ 都有一个表示
\[
 C(u)f(x)=G(u,g_i,h_i),
\]
其中 $C(u)$ 表示 $P[u]$ 中一个非零多项式，而 $G$ 在 $u,g_i,h_i$ 中为整。比较 $u$ 的某个适当幂的系数，就表明 $g_i(x)$ 与 $h_i(x)$ 给出 $\mathfrak S$ 的一个整基。于是\textbf{有限性判别准则已在一般情形中得到证明}。

\subsection*{\S\ 2. 不变量有限性}

1. 所谓\textbf{特征为 $p$ 的有限线性群}，是指由 $h$ 个线性替换 $A_1,\ldots,A_h$ 组成的一个群 $\mathfrak H$，这些替换的行列式非零，且系数来自特征为 $p$ 的一个域 $P$。这里 $A_k$ 表示线性替换
\[
\begin{aligned}
 x_1^{(k)}&=a_{11}^{(k)}x_1+\cdots+a_{1n}^{(k)}x_n,\\[-0.2em]
 &\ldots,\\[-0.2em]
 x_n^{(k)}&=a_{n1}^{(k)}x_1+\cdots+a_{nn}^{(k)}x_n,
\end{aligned}
\]
或者简记为 $(x^{(k)})=A_k(x)$。群 $\mathfrak H$ 因而把带有元素 $x_1,\ldots,x_n$ 的组 $(x)$ 变为带有元素 $x_1^{(k)},\ldots,x_n^{(k)}$ 的各组 $(x^{(k)})$。由于 $A_1,\ldots,A_h$ 中必须包含恒等变换，各组 $(x^{(k)})$ 中也包含组 $(x)$。

所谓该群的一个\textbf{整有理（绝对）不变量}，是指 $x_1,\ldots,x_n$ 的一个这样的整有理函数--系数来自 $P$\footnote{这并不限制一般性，因为任意一个系数来自某个特征为 $p$ 的域的不变量--为了定义这些不变量，该域必须与 $P$ 有一个共同的包域--都可以线性地分解为有限多个系数来自 $P$ 的不变量。一般情形中，只需把系数表示为关于 $P$ 线性无关的那些系数；于是，不变性就在这些独立系数中恒等地成立。}--它在应用 $A_1,\ldots,A_h$ 时恒等地保持不变。因此，各组 $(x^{(k)})$ 的所有对称函数--系数来自 $P$--都属于这些不变量。反过来，对每一个这样的不变量都有
\[
 f(x)=f(x^{(1)})=\cdots=f(x^{(h)})
 \quad\text{或}\quad
 h\cdot f(x)=f(x^{(1)})+\cdots+f(x^{(h)}).
\]
由此可知，在群的阶 $h$ 不被特征整除的情形--因而特别是在特征零情形--各组 $(x^{(k)})$ 的、系数来自 $P$ 的对称函数已经穷尽所有不变量。由于这些对称函数有一个有限整基，在这种情形下有限性证明就完成了；同时，整基给出为 Galois 预解式\footnote{见第28页脚注1。} 的系数系统 $G_{\alpha\alpha_1\ldots\alpha_n}(x)$：
\[
 \Phi(z,u)=\prod_{k=1}^{h}\bigl(z-u_1x_1^{(k)}-\cdots-u_nx_n^{(k)}\bigr)
\]
\[
 =z^h+\sum G_{\alpha\alpha_1\ldots\alpha_n}(x)
 z^{\alpha}u_1^{\alpha_1}\cdots u_n^{\alpha_n}
 \quad(\alpha\ne h;\ \alpha+\alpha_1+\cdots+\alpha_n=h).
\]

2. 在 $h$ 能被 $p$ 整除的情形，并不必然每个群不变量都是各组 $(x^{(k)})$ 的对称函数，因为现在不能再由 $h\cdot f(x)$ 推出 $f(x)$。但是，由有限多个不变量 $G_{\alpha\alpha_1\ldots\alpha_n}(x)$--即 Galois 预解式的系数--导出的环 $\mathfrak R$，构成所有不变量的系统 $\mathfrak S$ 的一个有限子环；还需证明，$\mathfrak S$ 关于 $\mathfrak R$ 满足有限性判别准则的前提。

首先，不失一般性\footnote{见第33页脚注。}，可把由全部替换系数 $a_{\mu\nu}^{(k)}$ 导出的 $P$ 的子域 $\overline P$ 作为系数域。这个由有限多个元素导出的域，是由向素域添加有限多个代数元素或超越元素而得到的；并且由于素域是完美的，遂有 $\overline P^{1/p}$ 关于 $\overline P$ 有限。\footnote{见第32页脚注2。}

还要证明 $\mathfrak S$ 在 $\mathfrak R$ 上代数地整依赖。按 1.，在 $h$ 个组 $x^{(k)}$ 中也出现原来的 $x_1,\ldots,x_n$；因此，当
\[
 z=u_1x_1+\cdots+u_nx_n
\]
时，Galois 预解式 $\Phi(z,u)$ 关于 $u$ 恒等为零。每次把除 $u_i$ 外的所有 $u$ 特化为零，并令 $u_i$ 等于单位元，就得到一个关系
\[
 x_i^h+\sum x_i^\alpha G_{\alpha0\ldots\alpha_i\ldots0}(x)=0
 \quad (\alpha\ne h;\ \alpha+\alpha_i=h),
\]
它表明每个 $x_i$ 都在 $\mathfrak R$ 上代数地整依赖。于是对于 $x$ 中任意多项式也同样成立；特别地，$\mathfrak S$ 在 $\mathfrak R$ 上代数地整依赖。因此有限性判别准则的前提满足，\textbf{不变量有限性得证}。

3. 还应指出，同样的推理给出\textbf{相对不变量与模不变量的有限性}。若一个多项式 $f(x)$ 的所有 $f(x^{(k)})$ 都只与 $f(x)$ 相差一个来自 $P$ 的非零因子，则称 $f(x)$ 为该群的一个\textbf{相对不变量}。绝对不变量的系统--特别是由 $G_{\alpha\alpha_1\ldots\alpha_n}(x)$ 导出的环 $\mathfrak R$--因而构成由所有相对不变量导出的环的一个有限子环；并且如上所述，关于 $\mathfrak R$ 的有限性判别准则前提满足。

若对于来自 $P$ 的每个特殊取值系统 $x_1=\alpha_1,\ldots,x_n=\alpha_n$，都有 $f(x)$ 等于 $f(x^{(k)})$，则称一个多项式 $f(x)$ 为该群的一个\textbf{模不变量}。每个绝对不变量同时也是模不变量；但如果 $P$ 只含有有限多个元素，反过来则未必成立。在这里，关于由 $G_{\alpha\alpha_1\ldots\alpha_n}(x)$ 导出的子环 $\mathfrak R$，有限性判别准则的前提同样满足。

\clearpage

% END live source-fidelity unit: translations/non_slavic/simplified_chinese/paper29/source_fidelity/v001/Noether_Paper29_SourceFidelity_SimplifiedChinese_v001.tex
\clearpage
% BEGIN live source-fidelity unit: translations/non_slavic/simplified_chinese/paper30/source_fidelity/v001/Noether_Paper30_SourceFidelity_SimplifiedChinese_v001.tex
\setcounter{footnote}{0}

% Source-fidelity v001: Paper30 complete Simplified Chinese translation unit.
